\section*{Table 1. Overview of the evidence base}

\begin{tabular}{p{0.24\textwidth} p{0.28\textwidth} p{0.40\textwidth}}
\toprule
Evidence domain & Main study designs & Overall pattern \\
\midrule
Occupational stress & Prospective cohort; analytical cross-sectional & Prospective studies linked high or increasing work stress with later MetS risk; cross-sectional findings were more mixed and often component-specific. \\
Diet and workplace nutrition & Analytical cross-sectional; randomised and quasi-experimental intervention & Healthier dietary patterns generally aligned with more favourable metabolic profiles, while workplace dietary interventions improved several MetS components. \\
Physical activity and sedentary behaviour & Cohort; analytical cross-sectional; randomised intervention & Leisure-time physical activity was generally associated with lower MetS risk, whereas prolonged sedentary exposure was associated with higher risk. \\
Multicomponent workplace interventions & Randomised and quasi-experimental intervention & Programmes combining diet, activity, counselling, self-monitoring or digital support generally improved anthropometric or metabolic outcomes. \\
\bottomrule
\end{tabular}

\vspace{0.5em}
\noindent\textit{Note:} Several studies contributed evidence to more than one domain. Detailed characteristics for all 34 studies are provided in Supplementary Table S1.
\section*{Table 2. Selected source-verified findings}

\begin{tabular}{p{0.19\textwidth} p{0.27\textwidth} p{0.45\textwidth}}
\toprule
Study & Exposure/intervention & Principal finding \\
\midrule
Garbarino \& Magnavita (2015) & Work stress & High work stress was associated with incident MetS (adjusted OR 2.68, 95\% CI 1.08--6.70). \\
Yamaguchi et al. (2018) & Change in job demands & Low-to-high job demands were associated with incident MetS (adjusted OR 3.27, 95\% CI 1.46--7.34). \\
Kuwahara et al. (2016) & Leisure-time exercise & Compared with the lowest category, mixed moderate/vigorous exercise was associated with lower MetS incidence (HR 0.74, 95\% CI 0.62--0.89 for 7.5--<16.5 MET-h/week). \\
Browne et al. (2017) & Physical activity & Active workers had lower odds of MetS than sedentary workers (OR 0.52, 95\% CI 0.27--0.98). \\
Motuma et al. (2020) & Sedentary behaviour and diet & Sedentary time $\geq$8 h/day was associated with higher MetS prevalence (APR 2.29, 95\% CI 1.88--2.80); high fruit/vegetable intake showed an inverse association. \\
Haufe et al. (2019) & Telemonitored exercise & Exercise reduced MetS severity versus control (between-group difference $-0.26$, 95\% CI $-0.35$ to $-0.16$). \\
Ma et al. (2025) & Canteen-based dietary intervention & The intervention improved fasting glucose, total and LDL cholesterol, waist circumference, BMI and HDL cholesterol. \\
Bogale et al. (2025) & Lifestyle education & Waist circumference decreased by 5.33 cm and systolic blood pressure by 6.96 mmHg, with additional lipid improvements. \\
Ayeltigah et al. (2026) & Work stress in military personnel & Work stress was associated with higher odds of MetS (OR 6.472, 95\% CI 1.535--27.292), although the estimate was imprecise. \\
Ryu et al. (2017) & Targeted workplace health promotion & Waist circumference and fasting glucose improved, while MetS prevalence declined from 41.5\% to 31.7\%. \\
\bottomrule
\end{tabular}

\vspace{0.5em}
\noindent\textit{Abbreviations:} APR, adjusted prevalence ratio; BMI, body mass index; CI, confidence interval; HR, hazard ratio; MetS, metabolic syndrome; OR, odds ratio.

\section*{Table 3. Methodological appraisal by design}

\begin{tabular}{p{0.24\textwidth} p{0.08\textwidth} p{0.24\textwidth} p{0.36\textwidth}}
\toprule
Design & n & Main strengths & Main limitations \\
\midrule
Randomised controlled trials & 6 & Stronger causal inference; standardised outcome assessment & Allocation concealment, masking, attrition and analysis by randomised group were incompletely reported in several trials. \\
Prospective/cohort studies & 6 & Temporal ordering; adjustment for major confounders; standardised MetS assessment & Incomplete reporting of attrition or missing-data handling in several studies. \\
Analytical cross-sectional studies & 17 & Broad occupational coverage; generally standardised MetS criteria & No temporal ordering; reverse causation remains plausible. \\
Quasi-/non-randomised interventions & 5 & Real-world workplace implementation & Non-randomised allocation and, in several studies, absence of a concurrent control group. \\
\bottomrule
\end{tabular}

\vspace{0.5em}
\noindent\textit{Note:} Design-specific Joanna Briggs Institute appraisal was used to qualify confidence in the narrative synthesis; no summed quality score was applied.